% Supplementary Methods - Content Only
% To be included via % Supplementary Methods - Content Only
% To be included via % Supplementary Methods - Content Only
% To be included via % Supplementary Methods - Content Only
% To be included via \input{Methods/supplementary_content}

% ========== Section 1: Confound Regression ==========

\subsection*{\textbf{S1. Confound Regression and Temporal Filtering}}

Within-run head motion was estimated via rigid-body realignment.
No temporal filtering or confound regression was applied; slow drift was modeled via linear per-run regressors in the general linear model.


% ========== Section 2: Quality Control ==========

\subsection*{\textbf{S2. Quality Control}}

Registration quality was evaluated by computing the intersection between the Wang Atlas ROIs (V1, V2, V3, hV4) and the BOLD brain mask in MNI space.
Across 10 subjects, mean ROI coverage was 84.3\% (SD = 21.7\%), and GLM valid ratio (percentage of ROI voxels with reliable stimulus-evoked responses) was 99.6\%.
We evaluated all subjects' data suitable for downstream analysis, though sub-07 showed reduced coverage (30.8\%) due to individual anatomical variability.


% ========== Section 3: Voxel Counts ==========



% ========== Section 5: Mean Activation ==========

\subsection*{\textbf{S5. Mean Activation Analysis}}

To verify that the shared-space geometric differences are not a signal-quality confound, we compared activation-level metrics between groups before SRM alignment. Mean activation magnitude (mean $|\beta|$), signal-to-noise ratio, color modulation depth, run-to-run reliability, and spatial variance did not differ between HC and CVD across V1, V2, V3, and hV4 (Welch's $t$-tests, all $p > 0.09$), and the two CVD cases fell within the HC range on every metric (Crawford--Howell, all $p > 0.09$). Voxel-level color selectivity was strong in both groups (one-way ANOVA across voxels, $F > 4$, $p < 0.001$ in V1--V3). Critically, none of the activation metrics correlated with SRM disparity across participants (Pearson $r$, 20 tests, all $p > 0.13$). The HC--CVD difference is therefore one of representational pattern, not signal amplitude: the cortical color response is equally strong in CVD but geometrically reorganized.


% ========== Section 6: Statistical Analysis ==========

\subsection*{\textbf{S6. Statistical Analysis}}

% [TODO: revise based on feedback] — FDR scope conflict: "All tests … FDR" overstates. FDR actually applies only to RDM color-pair comparisons (28 pairs) + the 6 identifiability checks; per-hue vulnerability and per-pair identification tests are uncorrected/exploratory. See PSYCH_THESIS/Methods/supplementary_content.tex for the revised wording.
All tests used a significance threshold of $\alpha = 0.05$ with False Discovery Rate correction \cite{benjamini1995}; $q = 0.05$ applied to RDM color-pair comparisons (28 pairs per participant).
Individual CVD participants were compared to the HC distribution using \citeA{crawford1998} modified t-tests:

\begin{equation}
t^* = \frac{x_{\text{CVD}} - \bar{x}_{\text{HC}}}{\text{SD}_{\text{HC}} \times \sqrt{(n+1)/n}}, \quad df = n-1 = 6
\end{equation}

with one-tailed p-values testing the directional hypothesis that CVD disparity exceeds HC.
Effect sizes were quantified via Hedges' g.
Group-level permutation tests (10,000 iterations) were used for HC-CVD comparisons by permuting subject labels; color-level permutation tests (1,000 iterations) were used for LOCO/LORO validation by permuting color labels within subjects.



% ========== Section 7: K-Selection ==========

\subsection*{\textbf{S7. Dimensionality Selection (K-Selection)}}

The reduced dimension $k$ for SRM was determined via leave-one-subject-out (LOSO) cross-validation based on within-subject representational dissimilarity matrix (RDM) reliability \cite{cunningham2014}.
For each candidate $k$ (range: 1--8) and fold, reliability was quantified as the Spearman correlation between RDMs computed from even runs (2,4,6) versus odd runs (1,3,5) of the held-out participant in SRM space.
We selected $k$ maximizing mean reliability across seven folds.
The selected values and their associated cross-subject RDM Spearman correlations were:

\begin{itemize}
\item V1: $k = 4$ (cross-subject RDM correlation = 0.60)
\item V2: $k = 4$ (cross-subject RDM correlation = 0.57)
\item V3: $k = 3$ (cross-subject RDM correlation = 0.55)
\item hV4: $k = 3$ (cross-subject RDM correlation = 0.32)
\end{itemize}


% ========== Section 8: GCV Details ==========

\subsection*{\textbf{S8. Generalized Cross-Validation for Ridge Regression}}
\label{app:gcv}

The regularization parameter $\alpha$ for ridge regression was selected via Generalized Cross-Validation (GCV; \citeNP{golub1979}), implemented as \texttt{sklearn.linear\_model.RidgeCV} with \texttt{gcv\_mode='auto'} (\citeNP{pedregosa2011}).
GCV provides a computationally efficient approximation to leave-one-out cross-validation for selecting $\alpha$.
The optimal $\alpha$ minimizes the GCV score:

\begin{equation}
\text{GCV}(\alpha) = \frac{1}{n} \|y - \hat{y}\|_2^2 / \left(1 - \frac{\text{tr}(H)}{n}\right)^2
\end{equation}

where $X$ denotes the design matrix, $y$ is the observed response vector, $\hat{y}$ is the ridge-predicted response, $H = X(X^T X + \alpha I)^{-1}X^T$ is the hat matrix, and $n$ is the number of training samples.


% ========== Section 9: LOO-Consistent Disparity ==========

\subsection*{\textbf{S9. Leave-One-Out Consistent Disparity Estimation}}

To quantify HC-CVD disparity with unbiased references, we used a leave-one-out (LOO) consistent procedure: for each fold $i$ (leaving out HC participant $i$), a reference pattern $R_{-i}$ was computed as the mean of the remaining six HC participants' aligned data.
Both the held-out $\text{HC}_i$ and both CVD participants were compared against this same reference $R_{-i}$, yielding disparity scores $d(\text{HC}_i, R_{-i})$ and $d(\text{CVD}_j, R_{-i})$.
This ensures HC and CVD disparities are measured against identical references within each fold.
Individual CVD disparity significance was assessed via \citeA{crawford1998} modified t-tests comparing each CVD participant's mean disparity (averaged across 7 folds) to the HC distribution ($n=7$, $df=6$).


% ========== Section 10: Cross-Validation Details ==========

\subsection*{\textbf{S10. Cross-Validation Procedures}}

\subsection*{\textbf{Leave-One-Run-Out (LORO)}}

LORO evaluated whether individual colors could be distinguished from voxel response amplitudes, testing the consistency of neural color encoding across runs.
For each of six runs, the corresponding response matrices ($B$, $C$) were held out as the test set, $W$ was trained on the remaining five runs.
Both color decoding and voxel response prediction were evaluated on the held-out run.

\subsection*{\textbf{Leave-One-Color-Out (LOCO)}}

LOCO tested whether geometric relationships of circular hue space were preserved in neural representation by evaluating interpolation to held-out colors.
For each color, all six run response matrices of the target color were held out. $W$ was trained on the remaining seven colors, and the model predicted either the held-out color from voxel responses or voxel responses from the held-out color's channel outputs.

\subsection*{\textbf{Leave-One-Subject-Out (LOSO)}}

LOSO evaluated zero-shot generalization to held-out participants.
As LOSO requires cross-subject prediction, SRM-transformed data were used to align voxel spaces across participants.
For each HC participant, $W$ was trained on the remaining HC participants' shared-space data and tested on the held-out participant.
This assessed whether the HC group-derived common space supported generalization to novel individuals.


% ========== Section 11: Evaluation Metrics ==========

\subsection*{\textbf{S11. Evaluation Metrics}}

\subsection*{\textbf{Color Decoding Metrics}}

Color decoding was evaluated using (1) mean absolute error (MAE) between predicted and true hue angles (chance = 90°, expected from a uniform random predictor on a circle), and (2) categorization accuracy with $\pm 22.5$° bins around each of the eight colors (chance = 12.5\% = 1/8).

\subsection*{\textbf{Voxel Prediction Metrics}}

Voxel response prediction was evaluated using Pearson correlation between predicted and observed response vectors, computed across all voxels and colors in the test set.

\subsection*{\textbf{Statistical Significance}}

Statistical significance of group comparisons (HC vs CVD) and individual CVD analyzes was assessed using nonparametric permutation tests (10,000 iterations; \citeNP{nichols2002}).


% ========== Section 12: SRM Residuals ==========

\subsection*{\textbf{S12. Shared Response Model Residuals}}

Each participant's residuals in the SRM model reflect the extent to which their data are not captured by the common structure, including inherent geometric differences.
The orthogonal constraint ($W_i^T W_i = I_k$) helps preserve geometric relationships between colors---distances and angles---in the common space.

Additionally, global between-subject RDM similarity was quantified using Spearman correlation across all 28 color pairs.


% ========== Section 13: Image Orientation ==========

\subsection*{\textbf{S13. Image Orientation Initialization}}

Image orientation was first initialized according to the scanner-defined obliquity information in the NIfTI header.
The mutual-information cost function was then optimized to maximize the statistical dependence between T1w and BOLD intensity distribution within the overlapping region.


% ========== Section 14: Procrustes Benefits ==========

\subsection*{\textbf{S14. Procrustes Alignment Benefits}}

This preprocessing reduces measurement noise across runs while preserving color representational geometry, thereby improving the quality of subsequent SRM fitting.


% ========== Section 16: Retinal-Family Model Comparison ==========

\subsection*{\textbf{S16. Comparison with Retinal-Family Distortion Models}}
\label{app:retinal_family}

Two retinal-family models were evaluated alongside the 2-component model by the same grid-search and permutation procedure (see Section~\ref{sec:methods:filter}).

\textbf{Models.} The 1-DOF Machado model~\citeNP{machado2009} shifts the L/M cone spectral sensitivity by $\Delta\lambda$~nm, yielding a hue-confusion pattern characteristic of the CVD subtype.
The 2-DOF retinal$+$cortical R$+$C model~\citeNP{tregillus2021} augments Machado with a cortical gain $g$: $r_{g}' = r_{g,\text{base}} + (1+g)(r_{g,\text{ret}} - r_{g,\text{base}})$, where $g = 0$ recovers the retinal-only model and $g = -1$ is exact cortical compensation.

\textbf{Fits (Table~\ref{tab:S1}).} Under the historical LOCO-$\rho$ argmax criterion (40{,}320 label permutations), reported here for reference only, retinal-family fits at hV4 yield Sub-08 R+C $(\Delta\lambda = 2.0$~nm, $g = +2.25$, $\rho = 0.857)$ and Sub-09 1-DOF Machado $(\Delta\lambda = 13.5$~nm, $\rho = 0.762)$; the latter shift magnitude falls within the established range for moderate--severe protanomaly~\cite{neitz2011}.
The large positive cortical gain for Sub-08 ($g = +2.25$, indicating $>$100\% overshoot beyond the exact Tregillus compensation level of $g = -1$;~\citeNP{tregillus2021}) is not physiologically interpretable as a gain parameter and is treated as an empirical characterization of residual hV4 structure.
The cross-subject model-class divergence (R+C for Sub-08; 1-DOF Machado for Sub-09) prevents a common-axis comparison of the two cases' distortion mechanisms.

\textbf{$\Delta$RDM cross-model convergence.} Using R+C parameters for Sub-09 ($\Delta\lambda = 19.5$~nm, $g = -1.10$), an independent $\Delta$RDM test reaches $p = 0.026$ (40{,}320 cosine-similarity permutations on crossnobis distances), providing cross-model corroboration for the Sub-09 distortion signal.
Retinal-family optimization directly against $\Delta$RDM fails for both CVD subjects; LOCO is the more discriminating phenotype and $\Delta$RDM serves as a convergence cross-check only.

\textbf{Arc-compression and LOCO vulnerability.} At Sub-09's severity, Machado arc-compression ($\Delta\lambda = 13.5$~nm) collapses three hues (c4 green, c5 cyan, c6 blue) to a single perceived angle in opponent hue space.
These three collapsed hues fall within Sub-09's LOCO vulnerability cluster (c5/c6 cyan/blue), confirming that arc-compression eliminates correction capacity precisely where Sub-09 requires it most.

\begin{table}[ht]
  \caption{Distortion-model fits at hV4 for both CVD participants.
  The 2-component model is listed at its adopted fit, selected by held-out composite test-loss ($\overline{L}_{\rm test}$, lower is better; main text, Section~\ref{sec:results:twocomp}).
  The retinal-family baselines were fit only under the historical LOCO-$\rho$ argmax criterion (Section~A.1 of Appendix~A; $\rho$ is Spearman correlation between simulated and observed per-color vulnerability), reported for reference only.}
  \label{tab:S1}
  \centering
  \small
  \begin{tabular}{llcccc}
    \toprule Subject & Model & DOF & Param.\ 1 & Param.\ 2 & Fit \\ \midrule Sub-08 (deutan) & 2-component & 2 & $\beta_s = 6^\circ$ & $\beta_c = -42^\circ$ & $\overline{L}_{\rm test} = -2.36$ \\ Sub-09 (protan) & 2-component & 2 & $\beta_s = 2^\circ$ & $\beta_c = +24^\circ$ & $\overline{L}_{\rm test} = -1.54$ \\ \midrule Sub-08 (deutan) & R+C\textsuperscript{a} & 2 & $\Delta\lambda = 2.0$~nm & $g = +2.25$ & $\rho = 0.857$ \\ Sub-09 (protan) & Machado\textsuperscript{a} & 1 & $\Delta\lambda = 13.5$~nm & --- & $\rho = 0.762$ \\ \bottomrule
  \end{tabular}
  \\[2pt]{\footnotesize \textsuperscript{a}Historical LOCO-$\rho$ argmax criterion (reference only); not the adopted filter fit.}
\end{table}

% ========== Section 17: HC magnitude-anchor distribution ==========

\subsection*{\textbf{S17. HC leave-one-out magnitude-anchor distribution}}
\label{app:hc_perm}

Because the production loss differs between participants (sub-08: $\gamma_{\rm OY} + L_{\rm RDM}^{(V2)}$; sub-09: $\gamma_{\rm all} + L_{\rm RDM}^{(V1)}$), the HC leave-one-out anchor was computed per participant under that participant's loss. Each of the $n = 7$ HC participants was refit as a pseudo-CVD case with the remaining HC pool as reference, yielding a distribution of fitted-parameter magnitudes $\|\hat\beta\|$: $30.5^\circ$ to $58.1^\circ$ (mean $49.1^\circ$) under the sub-08 loss, and $23.4^\circ$ to $55.5^\circ$ (mean $35.7^\circ$) under the sub-09 loss.
Both CVD estimates fall within their respective HC ranges (sub-08 $\|\hat\beta\| = 42.4^\circ$; sub-09 $\|\hat\beta\| = 24.1^\circ$), so parameter magnitude alone does not separate the CVD cases from HC; the distribution is reported as a descriptive magnitude anchor, not as a hypothesis test of CVD-versus-HC specificity.

% ========== Section 18: Effect Sizes ==========

\subsection*{\textbf{S18. Effect Sizes for Single-Case Comparisons}}
\label{app:effect_sizes}

Crawford--Howell single-case comparisons used the modified $t$ \citeNP{crawford1998}; the corresponding case-control effect size is $d_{cc} = (x_{\rm case} - \bar{x}_{\rm HC}) / s_{\rm HC} = t \cdot \sqrt{(n+1)/n}$ where $n$ is the HC sample size used in the test.
Mann--Whitney $U$ tests are reported with rank-biserial correlation $r_{\rm rb}$ as effect size.
Full per-test statistics ($t$, $U$, $p$, $d_{cc}$, $r_{\rm rb}$) for every Crawford--Howell and Mann--Whitney comparison reported in the main text are tabulated in the online repository (\texttt{statistics/effect\_sizes.csv}).

% NOTE: Bibliography는 main.tex에서 처리됨



% ========== Section 1: Confound Regression ==========

\subsection*{\textbf{S1. Confound Regression and Temporal Filtering}}

Within-run head motion was estimated via rigid-body realignment.
No temporal filtering or confound regression was applied; slow drift was modeled via linear per-run regressors in the general linear model.


% ========== Section 2: Quality Control ==========

\subsection*{\textbf{S2. Quality Control}}

Registration quality was evaluated by computing the intersection between the Wang Atlas ROIs (V1, V2, V3, hV4) and the BOLD brain mask in MNI space.
Across 10 subjects, mean ROI coverage was 84.3\% (SD = 21.7\%), and GLM valid ratio (percentage of ROI voxels with reliable stimulus-evoked responses) was 99.6\%.
We evaluated all subjects' data suitable for downstream analysis, though sub-07 showed reduced coverage (30.8\%) due to individual anatomical variability.


% ========== Section 3: Voxel Counts ==========



% ========== Section 5: Mean Activation ==========

\subsection*{\textbf{S5. Mean Activation Analysis}}

To verify that the shared-space geometric differences are not a signal-quality confound, we compared activation-level metrics between groups before SRM alignment. Mean activation magnitude (mean $|\beta|$), signal-to-noise ratio, color modulation depth, run-to-run reliability, and spatial variance did not differ between HC and CVD across V1, V2, V3, and hV4 (Welch's $t$-tests, all $p > 0.09$), and the two CVD cases fell within the HC range on every metric (Crawford--Howell, all $p > 0.09$). Voxel-level color selectivity was strong in both groups (one-way ANOVA across voxels, $F > 4$, $p < 0.001$ in V1--V3). Critically, none of the activation metrics correlated with SRM disparity across participants (Pearson $r$, 20 tests, all $p > 0.13$). The HC--CVD difference is therefore one of representational pattern, not signal amplitude: the cortical color response is equally strong in CVD but geometrically reorganized.


% ========== Section 6: Statistical Analysis ==========

\subsection*{\textbf{S6. Statistical Analysis}}

% [TODO: revise based on feedback] — FDR scope conflict: "All tests … FDR" overstates. FDR actually applies only to RDM color-pair comparisons (28 pairs) + the 6 identifiability checks; per-hue vulnerability and per-pair identification tests are uncorrected/exploratory. See PSYCH_THESIS/Methods/supplementary_content.tex for the revised wording.
All tests used a significance threshold of $\alpha = 0.05$ with False Discovery Rate correction \cite{benjamini1995}; $q = 0.05$ applied to RDM color-pair comparisons (28 pairs per participant).
Individual CVD participants were compared to the HC distribution using \citeA{crawford1998} modified t-tests:

\begin{equation}
t^* = \frac{x_{\text{CVD}} - \bar{x}_{\text{HC}}}{\text{SD}_{\text{HC}} \times \sqrt{(n+1)/n}}, \quad df = n-1 = 6
\end{equation}

with one-tailed p-values testing the directional hypothesis that CVD disparity exceeds HC.
Effect sizes were quantified via Hedges' g.
Group-level permutation tests (10,000 iterations) were used for HC-CVD comparisons by permuting subject labels; color-level permutation tests (1,000 iterations) were used for LOCO/LORO validation by permuting color labels within subjects.



% ========== Section 7: K-Selection ==========

\subsection*{\textbf{S7. Dimensionality Selection (K-Selection)}}

The reduced dimension $k$ for SRM was determined via leave-one-subject-out (LOSO) cross-validation based on within-subject representational dissimilarity matrix (RDM) reliability \cite{cunningham2014}.
For each candidate $k$ (range: 1--8) and fold, reliability was quantified as the Spearman correlation between RDMs computed from even runs (2,4,6) versus odd runs (1,3,5) of the held-out participant in SRM space.
We selected $k$ maximizing mean reliability across seven folds.
The selected values and their associated cross-subject RDM Spearman correlations were:

\begin{itemize}
\item V1: $k = 4$ (cross-subject RDM correlation = 0.60)
\item V2: $k = 4$ (cross-subject RDM correlation = 0.57)
\item V3: $k = 3$ (cross-subject RDM correlation = 0.55)
\item hV4: $k = 3$ (cross-subject RDM correlation = 0.32)
\end{itemize}


% ========== Section 8: GCV Details ==========

\subsection*{\textbf{S8. Generalized Cross-Validation for Ridge Regression}}
\label{app:gcv}

The regularization parameter $\alpha$ for ridge regression was selected via Generalized Cross-Validation (GCV; \citeNP{golub1979}), implemented as \texttt{sklearn.linear\_model.RidgeCV} with \texttt{gcv\_mode='auto'} (\citeNP{pedregosa2011}).
GCV provides a computationally efficient approximation to leave-one-out cross-validation for selecting $\alpha$.
The optimal $\alpha$ minimizes the GCV score:

\begin{equation}
\text{GCV}(\alpha) = \frac{1}{n} \|y - \hat{y}\|_2^2 / \left(1 - \frac{\text{tr}(H)}{n}\right)^2
\end{equation}

where $X$ denotes the design matrix, $y$ is the observed response vector, $\hat{y}$ is the ridge-predicted response, $H = X(X^T X + \alpha I)^{-1}X^T$ is the hat matrix, and $n$ is the number of training samples.


% ========== Section 9: LOO-Consistent Disparity ==========

\subsection*{\textbf{S9. Leave-One-Out Consistent Disparity Estimation}}

To quantify HC-CVD disparity with unbiased references, we used a leave-one-out (LOO) consistent procedure: for each fold $i$ (leaving out HC participant $i$), a reference pattern $R_{-i}$ was computed as the mean of the remaining six HC participants' aligned data.
Both the held-out $\text{HC}_i$ and both CVD participants were compared against this same reference $R_{-i}$, yielding disparity scores $d(\text{HC}_i, R_{-i})$ and $d(\text{CVD}_j, R_{-i})$.
This ensures HC and CVD disparities are measured against identical references within each fold.
Individual CVD disparity significance was assessed via \citeA{crawford1998} modified t-tests comparing each CVD participant's mean disparity (averaged across 7 folds) to the HC distribution ($n=7$, $df=6$).


% ========== Section 10: Cross-Validation Details ==========

\subsection*{\textbf{S10. Cross-Validation Procedures}}

\subsection*{\textbf{Leave-One-Run-Out (LORO)}}

LORO evaluated whether individual colors could be distinguished from voxel response amplitudes, testing the consistency of neural color encoding across runs.
For each of six runs, the corresponding response matrices ($B$, $C$) were held out as the test set, $W$ was trained on the remaining five runs.
Both color decoding and voxel response prediction were evaluated on the held-out run.

\subsection*{\textbf{Leave-One-Color-Out (LOCO)}}

LOCO tested whether geometric relationships of circular hue space were preserved in neural representation by evaluating interpolation to held-out colors.
For each color, all six run response matrices of the target color were held out. $W$ was trained on the remaining seven colors, and the model predicted either the held-out color from voxel responses or voxel responses from the held-out color's channel outputs.

\subsection*{\textbf{Leave-One-Subject-Out (LOSO)}}

LOSO evaluated zero-shot generalization to held-out participants.
As LOSO requires cross-subject prediction, SRM-transformed data were used to align voxel spaces across participants.
For each HC participant, $W$ was trained on the remaining HC participants' shared-space data and tested on the held-out participant.
This assessed whether the HC group-derived common space supported generalization to novel individuals.


% ========== Section 11: Evaluation Metrics ==========

\subsection*{\textbf{S11. Evaluation Metrics}}

\subsection*{\textbf{Color Decoding Metrics}}

Color decoding was evaluated using (1) mean absolute error (MAE) between predicted and true hue angles (chance = 90°, expected from a uniform random predictor on a circle), and (2) categorization accuracy with $\pm 22.5$° bins around each of the eight colors (chance = 12.5\% = 1/8).

\subsection*{\textbf{Voxel Prediction Metrics}}

Voxel response prediction was evaluated using Pearson correlation between predicted and observed response vectors, computed across all voxels and colors in the test set.

\subsection*{\textbf{Statistical Significance}}

Statistical significance of group comparisons (HC vs CVD) and individual CVD analyzes was assessed using nonparametric permutation tests (10,000 iterations; \citeNP{nichols2002}).


% ========== Section 12: SRM Residuals ==========

\subsection*{\textbf{S12. Shared Response Model Residuals}}

Each participant's residuals in the SRM model reflect the extent to which their data are not captured by the common structure, including inherent geometric differences.
The orthogonal constraint ($W_i^T W_i = I_k$) helps preserve geometric relationships between colors---distances and angles---in the common space.

Additionally, global between-subject RDM similarity was quantified using Spearman correlation across all 28 color pairs.


% ========== Section 13: Image Orientation ==========

\subsection*{\textbf{S13. Image Orientation Initialization}}

Image orientation was first initialized according to the scanner-defined obliquity information in the NIfTI header.
The mutual-information cost function was then optimized to maximize the statistical dependence between T1w and BOLD intensity distribution within the overlapping region.


% ========== Section 14: Procrustes Benefits ==========

\subsection*{\textbf{S14. Procrustes Alignment Benefits}}

This preprocessing reduces measurement noise across runs while preserving color representational geometry, thereby improving the quality of subsequent SRM fitting.


% ========== Section 16: Retinal-Family Model Comparison ==========

\subsection*{\textbf{S16. Comparison with Retinal-Family Distortion Models}}
\label{app:retinal_family}

Two retinal-family models were evaluated alongside the 2-component model by the same grid-search and permutation procedure (see Section~\ref{sec:methods:filter}).

\textbf{Models.} The 1-DOF Machado model~\citeNP{machado2009} shifts the L/M cone spectral sensitivity by $\Delta\lambda$~nm, yielding a hue-confusion pattern characteristic of the CVD subtype.
The 2-DOF retinal$+$cortical R$+$C model~\citeNP{tregillus2021} augments Machado with a cortical gain $g$: $r_{g}' = r_{g,\text{base}} + (1+g)(r_{g,\text{ret}} - r_{g,\text{base}})$, where $g = 0$ recovers the retinal-only model and $g = -1$ is exact cortical compensation.

\textbf{Fits (Table~\ref{tab:S1}).} Under the historical LOCO-$\rho$ argmax criterion (40{,}320 label permutations), reported here for reference only, retinal-family fits at hV4 yield Sub-08 R+C $(\Delta\lambda = 2.0$~nm, $g = +2.25$, $\rho = 0.857)$ and Sub-09 1-DOF Machado $(\Delta\lambda = 13.5$~nm, $\rho = 0.762)$; the latter shift magnitude falls within the established range for moderate--severe protanomaly~\cite{neitz2011}.
The large positive cortical gain for Sub-08 ($g = +2.25$, indicating $>$100\% overshoot beyond the exact Tregillus compensation level of $g = -1$;~\citeNP{tregillus2021}) is not physiologically interpretable as a gain parameter and is treated as an empirical characterization of residual hV4 structure.
The cross-subject model-class divergence (R+C for Sub-08; 1-DOF Machado for Sub-09) prevents a common-axis comparison of the two cases' distortion mechanisms.

\textbf{$\Delta$RDM cross-model convergence.} Using R+C parameters for Sub-09 ($\Delta\lambda = 19.5$~nm, $g = -1.10$), an independent $\Delta$RDM test reaches $p = 0.026$ (40{,}320 cosine-similarity permutations on crossnobis distances), providing cross-model corroboration for the Sub-09 distortion signal.
Retinal-family optimization directly against $\Delta$RDM fails for both CVD subjects; LOCO is the more discriminating phenotype and $\Delta$RDM serves as a convergence cross-check only.

\textbf{Arc-compression and LOCO vulnerability.} At Sub-09's severity, Machado arc-compression ($\Delta\lambda = 13.5$~nm) collapses three hues (c4 green, c5 cyan, c6 blue) to a single perceived angle in opponent hue space.
These three collapsed hues fall within Sub-09's LOCO vulnerability cluster (c5/c6 cyan/blue), confirming that arc-compression eliminates correction capacity precisely where Sub-09 requires it most.

\begin{table}[ht]
  \caption{Distortion-model fits at hV4 for both CVD participants.
  The 2-component model is listed at its adopted fit, selected by held-out composite test-loss ($\overline{L}_{\rm test}$, lower is better; main text, Section~\ref{sec:results:twocomp}).
  The retinal-family baselines were fit only under the historical LOCO-$\rho$ argmax criterion (Section~A.1 of Appendix~A; $\rho$ is Spearman correlation between simulated and observed per-color vulnerability), reported for reference only.}
  \label{tab:S1}
  \centering
  \small
  \begin{tabular}{llcccc}
    \toprule Subject & Model & DOF & Param.\ 1 & Param.\ 2 & Fit \\ \midrule Sub-08 (deutan) & 2-component & 2 & $\beta_s = 6^\circ$ & $\beta_c = -42^\circ$ & $\overline{L}_{\rm test} = -2.36$ \\ Sub-09 (protan) & 2-component & 2 & $\beta_s = 2^\circ$ & $\beta_c = +24^\circ$ & $\overline{L}_{\rm test} = -1.54$ \\ \midrule Sub-08 (deutan) & R+C\textsuperscript{a} & 2 & $\Delta\lambda = 2.0$~nm & $g = +2.25$ & $\rho = 0.857$ \\ Sub-09 (protan) & Machado\textsuperscript{a} & 1 & $\Delta\lambda = 13.5$~nm & --- & $\rho = 0.762$ \\ \bottomrule
  \end{tabular}
  \\[2pt]{\footnotesize \textsuperscript{a}Historical LOCO-$\rho$ argmax criterion (reference only); not the adopted filter fit.}
\end{table}

% ========== Section 17: HC magnitude-anchor distribution ==========

\subsection*{\textbf{S17. HC leave-one-out magnitude-anchor distribution}}
\label{app:hc_perm}

Because the production loss differs between participants (sub-08: $\gamma_{\rm OY} + L_{\rm RDM}^{(V2)}$; sub-09: $\gamma_{\rm all} + L_{\rm RDM}^{(V1)}$), the HC leave-one-out anchor was computed per participant under that participant's loss. Each of the $n = 7$ HC participants was refit as a pseudo-CVD case with the remaining HC pool as reference, yielding a distribution of fitted-parameter magnitudes $\|\hat\beta\|$: $30.5^\circ$ to $58.1^\circ$ (mean $49.1^\circ$) under the sub-08 loss, and $23.4^\circ$ to $55.5^\circ$ (mean $35.7^\circ$) under the sub-09 loss.
Both CVD estimates fall within their respective HC ranges (sub-08 $\|\hat\beta\| = 42.4^\circ$; sub-09 $\|\hat\beta\| = 24.1^\circ$), so parameter magnitude alone does not separate the CVD cases from HC; the distribution is reported as a descriptive magnitude anchor, not as a hypothesis test of CVD-versus-HC specificity.

% ========== Section 18: Effect Sizes ==========

\subsection*{\textbf{S18. Effect Sizes for Single-Case Comparisons}}
\label{app:effect_sizes}

Crawford--Howell single-case comparisons used the modified $t$ \citeNP{crawford1998}; the corresponding case-control effect size is $d_{cc} = (x_{\rm case} - \bar{x}_{\rm HC}) / s_{\rm HC} = t \cdot \sqrt{(n+1)/n}$ where $n$ is the HC sample size used in the test.
Mann--Whitney $U$ tests are reported with rank-biserial correlation $r_{\rm rb}$ as effect size.
Full per-test statistics ($t$, $U$, $p$, $d_{cc}$, $r_{\rm rb}$) for every Crawford--Howell and Mann--Whitney comparison reported in the main text are tabulated in the online repository (\texttt{statistics/effect\_sizes.csv}).

% NOTE: Bibliography는 main.tex에서 처리됨



% ========== Section 1: Confound Regression ==========

\subsection*{\textbf{S1. Confound Regression and Temporal Filtering}}

Within-run head motion was estimated via rigid-body realignment.
No temporal filtering or confound regression was applied; slow drift was modeled via linear per-run regressors in the general linear model.


% ========== Section 2: Quality Control ==========

\subsection*{\textbf{S2. Quality Control}}

Registration quality was evaluated by computing the intersection between the Wang Atlas ROIs (V1, V2, V3, hV4) and the BOLD brain mask in MNI space.
Across 10 subjects, mean ROI coverage was 84.3\% (SD = 21.7\%), and GLM valid ratio (percentage of ROI voxels with reliable stimulus-evoked responses) was 99.6\%.
We evaluated all subjects' data suitable for downstream analysis, though sub-07 showed reduced coverage (30.8\%) due to individual anatomical variability.


% ========== Section 3: Voxel Counts ==========



% ========== Section 5: Mean Activation ==========

\subsection*{\textbf{S5. Mean Activation Analysis}}

To verify that the shared-space geometric differences are not a signal-quality confound, we compared activation-level metrics between groups before SRM alignment. Mean activation magnitude (mean $|\beta|$), signal-to-noise ratio, color modulation depth, run-to-run reliability, and spatial variance did not differ between HC and CVD across V1, V2, V3, and hV4 (Welch's $t$-tests, all $p > 0.09$), and the two CVD cases fell within the HC range on every metric (Crawford--Howell, all $p > 0.09$). Voxel-level color selectivity was strong in both groups (one-way ANOVA across voxels, $F > 4$, $p < 0.001$ in V1--V3). Critically, none of the activation metrics correlated with SRM disparity across participants (Pearson $r$, 20 tests, all $p > 0.13$). The HC--CVD difference is therefore one of representational pattern, not signal amplitude: the cortical color response is equally strong in CVD but geometrically reorganized.


% ========== Section 6: Statistical Analysis ==========

\subsection*{\textbf{S6. Statistical Analysis}}

% [TODO: revise based on feedback] — FDR scope conflict: "All tests … FDR" overstates. FDR actually applies only to RDM color-pair comparisons (28 pairs) + the 6 identifiability checks; per-hue vulnerability and per-pair identification tests are uncorrected/exploratory. See PSYCH_THESIS/Methods/supplementary_content.tex for the revised wording.
All tests used a significance threshold of $\alpha = 0.05$ with False Discovery Rate correction \cite{benjamini1995}; $q = 0.05$ applied to RDM color-pair comparisons (28 pairs per participant).
Individual CVD participants were compared to the HC distribution using \citeA{crawford1998} modified t-tests:

\begin{equation}
t^* = \frac{x_{\text{CVD}} - \bar{x}_{\text{HC}}}{\text{SD}_{\text{HC}} \times \sqrt{(n+1)/n}}, \quad df = n-1 = 6
\end{equation}

with one-tailed p-values testing the directional hypothesis that CVD disparity exceeds HC.
Effect sizes were quantified via Hedges' g.
Group-level permutation tests (10,000 iterations) were used for HC-CVD comparisons by permuting subject labels; color-level permutation tests (1,000 iterations) were used for LOCO/LORO validation by permuting color labels within subjects.



% ========== Section 7: K-Selection ==========

\subsection*{\textbf{S7. Dimensionality Selection (K-Selection)}}

The reduced dimension $k$ for SRM was determined via leave-one-subject-out (LOSO) cross-validation based on within-subject representational dissimilarity matrix (RDM) reliability \cite{cunningham2014}.
For each candidate $k$ (range: 1--8) and fold, reliability was quantified as the Spearman correlation between RDMs computed from even runs (2,4,6) versus odd runs (1,3,5) of the held-out participant in SRM space.
We selected $k$ maximizing mean reliability across seven folds.
The selected values and their associated cross-subject RDM Spearman correlations were:

\begin{itemize}
\item V1: $k = 4$ (cross-subject RDM correlation = 0.60)
\item V2: $k = 4$ (cross-subject RDM correlation = 0.57)
\item V3: $k = 3$ (cross-subject RDM correlation = 0.55)
\item hV4: $k = 3$ (cross-subject RDM correlation = 0.32)
\end{itemize}


% ========== Section 8: GCV Details ==========

\subsection*{\textbf{S8. Generalized Cross-Validation for Ridge Regression}}
\label{app:gcv}

The regularization parameter $\alpha$ for ridge regression was selected via Generalized Cross-Validation (GCV; \citeNP{golub1979}), implemented as \texttt{sklearn.linear\_model.RidgeCV} with \texttt{gcv\_mode='auto'} (\citeNP{pedregosa2011}).
GCV provides a computationally efficient approximation to leave-one-out cross-validation for selecting $\alpha$.
The optimal $\alpha$ minimizes the GCV score:

\begin{equation}
\text{GCV}(\alpha) = \frac{1}{n} \|y - \hat{y}\|_2^2 / \left(1 - \frac{\text{tr}(H)}{n}\right)^2
\end{equation}

where $X$ denotes the design matrix, $y$ is the observed response vector, $\hat{y}$ is the ridge-predicted response, $H = X(X^T X + \alpha I)^{-1}X^T$ is the hat matrix, and $n$ is the number of training samples.


% ========== Section 9: LOO-Consistent Disparity ==========

\subsection*{\textbf{S9. Leave-One-Out Consistent Disparity Estimation}}

To quantify HC-CVD disparity with unbiased references, we used a leave-one-out (LOO) consistent procedure: for each fold $i$ (leaving out HC participant $i$), a reference pattern $R_{-i}$ was computed as the mean of the remaining six HC participants' aligned data.
Both the held-out $\text{HC}_i$ and both CVD participants were compared against this same reference $R_{-i}$, yielding disparity scores $d(\text{HC}_i, R_{-i})$ and $d(\text{CVD}_j, R_{-i})$.
This ensures HC and CVD disparities are measured against identical references within each fold.
Individual CVD disparity significance was assessed via \citeA{crawford1998} modified t-tests comparing each CVD participant's mean disparity (averaged across 7 folds) to the HC distribution ($n=7$, $df=6$).


% ========== Section 10: Cross-Validation Details ==========

\subsection*{\textbf{S10. Cross-Validation Procedures}}

\subsection*{\textbf{Leave-One-Run-Out (LORO)}}

LORO evaluated whether individual colors could be distinguished from voxel response amplitudes, testing the consistency of neural color encoding across runs.
For each of six runs, the corresponding response matrices ($B$, $C$) were held out as the test set, $W$ was trained on the remaining five runs.
Both color decoding and voxel response prediction were evaluated on the held-out run.

\subsection*{\textbf{Leave-One-Color-Out (LOCO)}}

LOCO tested whether geometric relationships of circular hue space were preserved in neural representation by evaluating interpolation to held-out colors.
For each color, all six run response matrices of the target color were held out. $W$ was trained on the remaining seven colors, and the model predicted either the held-out color from voxel responses or voxel responses from the held-out color's channel outputs.

\subsection*{\textbf{Leave-One-Subject-Out (LOSO)}}

LOSO evaluated zero-shot generalization to held-out participants.
As LOSO requires cross-subject prediction, SRM-transformed data were used to align voxel spaces across participants.
For each HC participant, $W$ was trained on the remaining HC participants' shared-space data and tested on the held-out participant.
This assessed whether the HC group-derived common space supported generalization to novel individuals.


% ========== Section 11: Evaluation Metrics ==========

\subsection*{\textbf{S11. Evaluation Metrics}}

\subsection*{\textbf{Color Decoding Metrics}}

Color decoding was evaluated using (1) mean absolute error (MAE) between predicted and true hue angles (chance = 90°, expected from a uniform random predictor on a circle), and (2) categorization accuracy with $\pm 22.5$° bins around each of the eight colors (chance = 12.5\% = 1/8).

\subsection*{\textbf{Voxel Prediction Metrics}}

Voxel response prediction was evaluated using Pearson correlation between predicted and observed response vectors, computed across all voxels and colors in the test set.

\subsection*{\textbf{Statistical Significance}}

Statistical significance of group comparisons (HC vs CVD) and individual CVD analyzes was assessed using nonparametric permutation tests (10,000 iterations; \citeNP{nichols2002}).


% ========== Section 12: SRM Residuals ==========

\subsection*{\textbf{S12. Shared Response Model Residuals}}

Each participant's residuals in the SRM model reflect the extent to which their data are not captured by the common structure, including inherent geometric differences.
The orthogonal constraint ($W_i^T W_i = I_k$) helps preserve geometric relationships between colors---distances and angles---in the common space.

Additionally, global between-subject RDM similarity was quantified using Spearman correlation across all 28 color pairs.


% ========== Section 13: Image Orientation ==========

\subsection*{\textbf{S13. Image Orientation Initialization}}

Image orientation was first initialized according to the scanner-defined obliquity information in the NIfTI header.
The mutual-information cost function was then optimized to maximize the statistical dependence between T1w and BOLD intensity distribution within the overlapping region.


% ========== Section 14: Procrustes Benefits ==========

\subsection*{\textbf{S14. Procrustes Alignment Benefits}}

This preprocessing reduces measurement noise across runs while preserving color representational geometry, thereby improving the quality of subsequent SRM fitting.


% ========== Section 16: Retinal-Family Model Comparison ==========

\subsection*{\textbf{S16. Comparison with Retinal-Family Distortion Models}}
\label{app:retinal_family}

Two retinal-family models were evaluated alongside the 2-component model by the same grid-search and permutation procedure (see Section~\ref{sec:methods:filter}).

\textbf{Models.} The 1-DOF Machado model~\citeNP{machado2009} shifts the L/M cone spectral sensitivity by $\Delta\lambda$~nm, yielding a hue-confusion pattern characteristic of the CVD subtype.
The 2-DOF retinal$+$cortical R$+$C model~\citeNP{tregillus2021} augments Machado with a cortical gain $g$: $r_{g}' = r_{g,\text{base}} + (1+g)(r_{g,\text{ret}} - r_{g,\text{base}})$, where $g = 0$ recovers the retinal-only model and $g = -1$ is exact cortical compensation.

\textbf{Fits (Table~\ref{tab:S1}).} Under the historical LOCO-$\rho$ argmax criterion (40{,}320 label permutations), reported here for reference only, retinal-family fits at hV4 yield Sub-08 R+C $(\Delta\lambda = 2.0$~nm, $g = +2.25$, $\rho = 0.857)$ and Sub-09 1-DOF Machado $(\Delta\lambda = 13.5$~nm, $\rho = 0.762)$; the latter shift magnitude falls within the established range for moderate--severe protanomaly~\cite{neitz2011}.
The large positive cortical gain for Sub-08 ($g = +2.25$, indicating $>$100\% overshoot beyond the exact Tregillus compensation level of $g = -1$;~\citeNP{tregillus2021}) is not physiologically interpretable as a gain parameter and is treated as an empirical characterization of residual hV4 structure.
The cross-subject model-class divergence (R+C for Sub-08; 1-DOF Machado for Sub-09) prevents a common-axis comparison of the two cases' distortion mechanisms.

\textbf{$\Delta$RDM cross-model convergence.} Using R+C parameters for Sub-09 ($\Delta\lambda = 19.5$~nm, $g = -1.10$), an independent $\Delta$RDM test reaches $p = 0.026$ (40{,}320 cosine-similarity permutations on crossnobis distances), providing cross-model corroboration for the Sub-09 distortion signal.
Retinal-family optimization directly against $\Delta$RDM fails for both CVD subjects; LOCO is the more discriminating phenotype and $\Delta$RDM serves as a convergence cross-check only.

\textbf{Arc-compression and LOCO vulnerability.} At Sub-09's severity, Machado arc-compression ($\Delta\lambda = 13.5$~nm) collapses three hues (c4 green, c5 cyan, c6 blue) to a single perceived angle in opponent hue space.
These three collapsed hues fall within Sub-09's LOCO vulnerability cluster (c5/c6 cyan/blue), confirming that arc-compression eliminates correction capacity precisely where Sub-09 requires it most.

\begin{table}[ht]
  \caption{Distortion-model fits at hV4 for both CVD participants.
  The 2-component model is listed at its adopted fit, selected by held-out composite test-loss ($\overline{L}_{\rm test}$, lower is better; main text, Section~\ref{sec:results:twocomp}).
  The retinal-family baselines were fit only under the historical LOCO-$\rho$ argmax criterion (Section~A.1 of Appendix~A; $\rho$ is Spearman correlation between simulated and observed per-color vulnerability), reported for reference only.}
  \label{tab:S1}
  \centering
  \small
  \begin{tabular}{llcccc}
    \toprule Subject & Model & DOF & Param.\ 1 & Param.\ 2 & Fit \\ \midrule Sub-08 (deutan) & 2-component & 2 & $\beta_s = 6^\circ$ & $\beta_c = -42^\circ$ & $\overline{L}_{\rm test} = -2.36$ \\ Sub-09 (protan) & 2-component & 2 & $\beta_s = 2^\circ$ & $\beta_c = +24^\circ$ & $\overline{L}_{\rm test} = -1.54$ \\ \midrule Sub-08 (deutan) & R+C\textsuperscript{a} & 2 & $\Delta\lambda = 2.0$~nm & $g = +2.25$ & $\rho = 0.857$ \\ Sub-09 (protan) & Machado\textsuperscript{a} & 1 & $\Delta\lambda = 13.5$~nm & --- & $\rho = 0.762$ \\ \bottomrule
  \end{tabular}
  \\[2pt]{\footnotesize \textsuperscript{a}Historical LOCO-$\rho$ argmax criterion (reference only); not the adopted filter fit.}
\end{table}

% ========== Section 17: HC magnitude-anchor distribution ==========

\subsection*{\textbf{S17. HC leave-one-out magnitude-anchor distribution}}
\label{app:hc_perm}

Because the production loss differs between participants (sub-08: $\gamma_{\rm OY} + L_{\rm RDM}^{(V2)}$; sub-09: $\gamma_{\rm all} + L_{\rm RDM}^{(V1)}$), the HC leave-one-out anchor was computed per participant under that participant's loss. Each of the $n = 7$ HC participants was refit as a pseudo-CVD case with the remaining HC pool as reference, yielding a distribution of fitted-parameter magnitudes $\|\hat\beta\|$: $30.5^\circ$ to $58.1^\circ$ (mean $49.1^\circ$) under the sub-08 loss, and $23.4^\circ$ to $55.5^\circ$ (mean $35.7^\circ$) under the sub-09 loss.
Both CVD estimates fall within their respective HC ranges (sub-08 $\|\hat\beta\| = 42.4^\circ$; sub-09 $\|\hat\beta\| = 24.1^\circ$), so parameter magnitude alone does not separate the CVD cases from HC; the distribution is reported as a descriptive magnitude anchor, not as a hypothesis test of CVD-versus-HC specificity.

% ========== Section 18: Effect Sizes ==========

\subsection*{\textbf{S18. Effect Sizes for Single-Case Comparisons}}
\label{app:effect_sizes}

Crawford--Howell single-case comparisons used the modified $t$ \citeNP{crawford1998}; the corresponding case-control effect size is $d_{cc} = (x_{\rm case} - \bar{x}_{\rm HC}) / s_{\rm HC} = t \cdot \sqrt{(n+1)/n}$ where $n$ is the HC sample size used in the test.
Mann--Whitney $U$ tests are reported with rank-biserial correlation $r_{\rm rb}$ as effect size.
Full per-test statistics ($t$, $U$, $p$, $d_{cc}$, $r_{\rm rb}$) for every Crawford--Howell and Mann--Whitney comparison reported in the main text are tabulated in the online repository (\texttt{statistics/effect\_sizes.csv}).

% NOTE: Bibliography는 main.tex에서 처리됨



% ========== Section 1: Confound Regression ==========

\subsection*{\textbf{S1. Confound Regression and Temporal Filtering}}

Within-run head motion was estimated via rigid-body realignment.
No temporal filtering or confound regression was applied; slow drift was modeled via linear per-run regressors in the general linear model.


% ========== Section 2: Quality Control ==========

\subsection*{\textbf{S2. Quality Control}}

Registration quality was evaluated by computing the intersection between the Wang Atlas ROIs (V1, V2, V3, hV4) and the BOLD brain mask in MNI space.
Across 10 subjects, mean ROI coverage was 84.3\% (SD = 21.7\%), and GLM valid ratio (percentage of ROI voxels with reliable stimulus-evoked responses) was 99.6\%.
We evaluated all subjects' data suitable for downstream analysis, though sub-07 showed reduced coverage (30.8\%) due to individual anatomical variability.


% ========== Section 3: Voxel Counts ==========



% ========== Section 5: Mean Activation ==========

\subsection*{\textbf{S5. Mean Activation Analysis}}

To verify that the shared-space geometric differences are not a signal-quality confound, we compared activation-level metrics between groups before SRM alignment. Mean activation magnitude (mean $|\beta|$), signal-to-noise ratio, color modulation depth, run-to-run reliability, and spatial variance did not differ between HC and CVD across V1, V2, V3, and hV4 (Welch's $t$-tests, all $p > 0.09$), and the two CVD cases fell within the HC range on every metric (Crawford--Howell, all $p > 0.09$). Voxel-level color selectivity was strong in both groups (one-way ANOVA across voxels, $F > 4$, $p < 0.001$ in V1--V3). Critically, none of the activation metrics correlated with SRM disparity across participants (Pearson $r$, 20 tests, all $p > 0.13$). The HC--CVD difference is therefore one of representational pattern, not signal amplitude: the cortical color response is equally strong in CVD but geometrically reorganized.


% ========== Section 6: Statistical Analysis ==========

\subsection*{\textbf{S6. Statistical Analysis}}

% [TODO: revise based on feedback] — FDR scope conflict: "All tests … FDR" overstates. FDR actually applies only to RDM color-pair comparisons (28 pairs) + the 6 identifiability checks; per-hue vulnerability and per-pair identification tests are uncorrected/exploratory. See PSYCH_THESIS/Methods/supplementary_content.tex for the revised wording.
All tests used a significance threshold of $\alpha = 0.05$ with False Discovery Rate correction \cite{benjamini1995}; $q = 0.05$ applied to RDM color-pair comparisons (28 pairs per participant).
Individual CVD participants were compared to the HC distribution using \citeA{crawford1998} modified t-tests:

\begin{equation}
t^* = \frac{x_{\text{CVD}} - \bar{x}_{\text{HC}}}{\text{SD}_{\text{HC}} \times \sqrt{(n+1)/n}}, \quad df = n-1 = 6
\end{equation}

with one-tailed p-values testing the directional hypothesis that CVD disparity exceeds HC.
Effect sizes were quantified via Hedges' g.
Group-level permutation tests (10,000 iterations) were used for HC-CVD comparisons by permuting subject labels; color-level permutation tests (1,000 iterations) were used for LOCO/LORO validation by permuting color labels within subjects.



% ========== Section 7: K-Selection ==========

\subsection*{\textbf{S7. Dimensionality Selection (K-Selection)}}

The reduced dimension $k$ for SRM was determined via leave-one-subject-out (LOSO) cross-validation based on within-subject representational dissimilarity matrix (RDM) reliability \cite{cunningham2014}.
For each candidate $k$ (range: 1--8) and fold, reliability was quantified as the Spearman correlation between RDMs computed from even runs (2,4,6) versus odd runs (1,3,5) of the held-out participant in SRM space.
We selected $k$ maximizing mean reliability across seven folds.
The selected values and their associated cross-subject RDM Spearman correlations were:

\begin{itemize}
\item V1: $k = 4$ (cross-subject RDM correlation = 0.60)
\item V2: $k = 4$ (cross-subject RDM correlation = 0.57)
\item V3: $k = 3$ (cross-subject RDM correlation = 0.55)
\item hV4: $k = 3$ (cross-subject RDM correlation = 0.32)
\end{itemize}


% ========== Section 8: GCV Details ==========

\subsection*{\textbf{S8. Generalized Cross-Validation for Ridge Regression}}
\label{app:gcv}

The regularization parameter $\alpha$ for ridge regression was selected via Generalized Cross-Validation (GCV; \citeNP{golub1979}), implemented as \texttt{sklearn.linear\_model.RidgeCV} with \texttt{gcv\_mode='auto'} (\citeNP{pedregosa2011}).
GCV provides a computationally efficient approximation to leave-one-out cross-validation for selecting $\alpha$.
The optimal $\alpha$ minimizes the GCV score:

\begin{equation}
\text{GCV}(\alpha) = \frac{1}{n} \|y - \hat{y}\|_2^2 / \left(1 - \frac{\text{tr}(H)}{n}\right)^2
\end{equation}

where $X$ denotes the design matrix, $y$ is the observed response vector, $\hat{y}$ is the ridge-predicted response, $H = X(X^T X + \alpha I)^{-1}X^T$ is the hat matrix, and $n$ is the number of training samples.


% ========== Section 9: LOO-Consistent Disparity ==========

\subsection*{\textbf{S9. Leave-One-Out Consistent Disparity Estimation}}

To quantify HC-CVD disparity with unbiased references, we used a leave-one-out (LOO) consistent procedure: for each fold $i$ (leaving out HC participant $i$), a reference pattern $R_{-i}$ was computed as the mean of the remaining six HC participants' aligned data.
Both the held-out $\text{HC}_i$ and both CVD participants were compared against this same reference $R_{-i}$, yielding disparity scores $d(\text{HC}_i, R_{-i})$ and $d(\text{CVD}_j, R_{-i})$.
This ensures HC and CVD disparities are measured against identical references within each fold.
Individual CVD disparity significance was assessed via \citeA{crawford1998} modified t-tests comparing each CVD participant's mean disparity (averaged across 7 folds) to the HC distribution ($n=7$, $df=6$).


% ========== Section 10: Cross-Validation Details ==========

\subsection*{\textbf{S10. Cross-Validation Procedures}}

\subsection*{\textbf{Leave-One-Run-Out (LORO)}}

LORO evaluated whether individual colors could be distinguished from voxel response amplitudes, testing the consistency of neural color encoding across runs.
For each of six runs, the corresponding response matrices ($B$, $C$) were held out as the test set, $W$ was trained on the remaining five runs.
Both color decoding and voxel response prediction were evaluated on the held-out run.

\subsection*{\textbf{Leave-One-Color-Out (LOCO)}}

LOCO tested whether geometric relationships of circular hue space were preserved in neural representation by evaluating interpolation to held-out colors.
For each color, all six run response matrices of the target color were held out. $W$ was trained on the remaining seven colors, and the model predicted either the held-out color from voxel responses or voxel responses from the held-out color's channel outputs.

\subsection*{\textbf{Leave-One-Subject-Out (LOSO)}}

LOSO evaluated zero-shot generalization to held-out participants.
As LOSO requires cross-subject prediction, SRM-transformed data were used to align voxel spaces across participants.
For each HC participant, $W$ was trained on the remaining HC participants' shared-space data and tested on the held-out participant.
This assessed whether the HC group-derived common space supported generalization to novel individuals.


% ========== Section 11: Evaluation Metrics ==========

\subsection*{\textbf{S11. Evaluation Metrics}}

\subsection*{\textbf{Color Decoding Metrics}}

Color decoding was evaluated using (1) mean absolute error (MAE) between predicted and true hue angles (chance = 90°, expected from a uniform random predictor on a circle), and (2) categorization accuracy with $\pm 22.5$° bins around each of the eight colors (chance = 12.5\% = 1/8).

\subsection*{\textbf{Voxel Prediction Metrics}}

Voxel response prediction was evaluated using Pearson correlation between predicted and observed response vectors, computed across all voxels and colors in the test set.

\subsection*{\textbf{Statistical Significance}}

Statistical significance of group comparisons (HC vs CVD) and individual CVD analyzes was assessed using nonparametric permutation tests (10,000 iterations; \citeNP{nichols2002}).


% ========== Section 12: SRM Residuals ==========

\subsection*{\textbf{S12. Shared Response Model Residuals}}

Each participant's residuals in the SRM model reflect the extent to which their data are not captured by the common structure, including inherent geometric differences.
The orthogonal constraint ($W_i^T W_i = I_k$) helps preserve geometric relationships between colors---distances and angles---in the common space.

Additionally, global between-subject RDM similarity was quantified using Spearman correlation across all 28 color pairs.


% ========== Section 13: Image Orientation ==========

\subsection*{\textbf{S13. Image Orientation Initialization}}

Image orientation was first initialized according to the scanner-defined obliquity information in the NIfTI header.
The mutual-information cost function was then optimized to maximize the statistical dependence between T1w and BOLD intensity distribution within the overlapping region.


% ========== Section 14: Procrustes Benefits ==========

\subsection*{\textbf{S14. Procrustes Alignment Benefits}}

This preprocessing reduces measurement noise across runs while preserving color representational geometry, thereby improving the quality of subsequent SRM fitting.


% ========== Section 16: Retinal-Family Model Comparison ==========

\subsection*{\textbf{S16. Comparison with Retinal-Family Distortion Models}}
\label{app:retinal_family}

Two retinal-family models were evaluated alongside the 2-component model by the same grid-search and permutation procedure (see Section~\ref{sec:methods:filter}).

\textbf{Models.} The 1-DOF Machado model~\citeNP{machado2009} shifts the L/M cone spectral sensitivity by $\Delta\lambda$~nm, yielding a hue-confusion pattern characteristic of the CVD subtype.
The 2-DOF retinal$+$cortical R$+$C model~\citeNP{tregillus2021} augments Machado with a cortical gain $g$: $r_{g}' = r_{g,\text{base}} + (1+g)(r_{g,\text{ret}} - r_{g,\text{base}})$, where $g = 0$ recovers the retinal-only model and $g = -1$ is exact cortical compensation.

\textbf{Fits (Table~\ref{tab:S1}).} Under the historical LOCO-$\rho$ argmax criterion (40{,}320 label permutations), reported here for reference only, retinal-family fits at hV4 yield Sub-08 R+C $(\Delta\lambda = 2.0$~nm, $g = +2.25$, $\rho = 0.857)$ and Sub-09 1-DOF Machado $(\Delta\lambda = 13.5$~nm, $\rho = 0.762)$; the latter shift magnitude falls within the established range for moderate--severe protanomaly~\cite{neitz2011}.
The large positive cortical gain for Sub-08 ($g = +2.25$, indicating $>$100\% overshoot beyond the exact Tregillus compensation level of $g = -1$;~\citeNP{tregillus2021}) is not physiologically interpretable as a gain parameter and is treated as an empirical characterization of residual hV4 structure.
The cross-subject model-class divergence (R+C for Sub-08; 1-DOF Machado for Sub-09) prevents a common-axis comparison of the two cases' distortion mechanisms.

\textbf{$\Delta$RDM cross-model convergence.} Using R+C parameters for Sub-09 ($\Delta\lambda = 19.5$~nm, $g = -1.10$), an independent $\Delta$RDM test reaches $p = 0.026$ (40{,}320 cosine-similarity permutations on crossnobis distances), providing cross-model corroboration for the Sub-09 distortion signal.
Retinal-family optimization directly against $\Delta$RDM fails for both CVD subjects; LOCO is the more discriminating phenotype and $\Delta$RDM serves as a convergence cross-check only.

\textbf{Arc-compression and LOCO vulnerability.} At Sub-09's severity, Machado arc-compression ($\Delta\lambda = 13.5$~nm) collapses three hues (c4 green, c5 cyan, c6 blue) to a single perceived angle in opponent hue space.
These three collapsed hues fall within Sub-09's LOCO vulnerability cluster (c5/c6 cyan/blue), confirming that arc-compression eliminates correction capacity precisely where Sub-09 requires it most.

\begin{table}[ht]
  \caption{Distortion-model fits at hV4 for both CVD participants.
  The 2-component model is listed at its adopted fit, selected by held-out composite test-loss ($\overline{L}_{\rm test}$, lower is better; main text, Section~\ref{sec:results:twocomp}).
  The retinal-family baselines were fit only under the historical LOCO-$\rho$ argmax criterion (Section~A.1 of Appendix~A; $\rho$ is Spearman correlation between simulated and observed per-color vulnerability), reported for reference only.}
  \label{tab:S1}
  \centering
  \small
  \begin{tabular}{llcccc}
    \toprule Subject & Model & DOF & Param.\ 1 & Param.\ 2 & Fit \\ \midrule Sub-08 (deutan) & 2-component & 2 & $\beta_s = 6^\circ$ & $\beta_c = -42^\circ$ & $\overline{L}_{\rm test} = -2.36$ \\ Sub-09 (protan) & 2-component & 2 & $\beta_s = 2^\circ$ & $\beta_c = +24^\circ$ & $\overline{L}_{\rm test} = -1.54$ \\ \midrule Sub-08 (deutan) & R+C\textsuperscript{a} & 2 & $\Delta\lambda = 2.0$~nm & $g = +2.25$ & $\rho = 0.857$ \\ Sub-09 (protan) & Machado\textsuperscript{a} & 1 & $\Delta\lambda = 13.5$~nm & --- & $\rho = 0.762$ \\ \bottomrule
  \end{tabular}
  \\[2pt]{\footnotesize \textsuperscript{a}Historical LOCO-$\rho$ argmax criterion (reference only); not the adopted filter fit.}
\end{table}

% ========== Section 17: HC magnitude-anchor distribution ==========

\subsection*{\textbf{S17. HC leave-one-out magnitude-anchor distribution}}
\label{app:hc_perm}

Because the production loss differs between participants (sub-08: $\gamma_{\rm OY} + L_{\rm RDM}^{(V2)}$; sub-09: $\gamma_{\rm all} + L_{\rm RDM}^{(V1)}$), the HC leave-one-out anchor was computed per participant under that participant's loss. Each of the $n = 7$ HC participants was refit as a pseudo-CVD case with the remaining HC pool as reference, yielding a distribution of fitted-parameter magnitudes $\|\hat\beta\|$: $30.5^\circ$ to $58.1^\circ$ (mean $49.1^\circ$) under the sub-08 loss, and $23.4^\circ$ to $55.5^\circ$ (mean $35.7^\circ$) under the sub-09 loss.
Both CVD estimates fall within their respective HC ranges (sub-08 $\|\hat\beta\| = 42.4^\circ$; sub-09 $\|\hat\beta\| = 24.1^\circ$), so parameter magnitude alone does not separate the CVD cases from HC; the distribution is reported as a descriptive magnitude anchor, not as a hypothesis test of CVD-versus-HC specificity.

% ========== Section 18: Effect Sizes ==========

\subsection*{\textbf{S18. Effect Sizes for Single-Case Comparisons}}
\label{app:effect_sizes}

Crawford--Howell single-case comparisons used the modified $t$ \citeNP{crawford1998}; the corresponding case-control effect size is $d_{cc} = (x_{\rm case} - \bar{x}_{\rm HC}) / s_{\rm HC} = t \cdot \sqrt{(n+1)/n}$ where $n$ is the HC sample size used in the test.
Mann--Whitney $U$ tests are reported with rank-biserial correlation $r_{\rm rb}$ as effect size.
Full per-test statistics ($t$, $U$, $p$, $d_{cc}$, $r_{\rm rb}$) for every Crawford--Howell and Mann--Whitney comparison reported in the main text are tabulated in the online repository (\texttt{statistics/effect\_sizes.csv}).

% NOTE: Bibliography는 main.tex에서 처리됨

